\section{Discussion}

\subsection{Interpretation: The Self-Reading Mechanism}

Our results support a ``self-reading'' interpretation of how CoT+confidence may improve calibration. When prompted to reason step-by-step, the model generates linguistic markers of epistemic uncertainty---hedging words like ``may,'' ``could,'' and ``however.'' Due to autoregressive generation, these markers remain in the model's context when it subsequently generates a confidence estimate. The significant negative correlation ($r=-0.315$) between hedging count and confidence indicates that the model adjusts its confidence judgment downward when more uncertainty signals are present in its own reasoning.

\subsection{Alternative Explanations}

\textbf{Difficulty Confound:} One alternative is that difficult questions produce both more hedging and lower confidence, without the hedging causally affecting confidence. This interpretation still supports using hedging as a calibration signal.

\textbf{Prompt Artifact:} Another alternative is that our prompt structure induces the correlation through some artifact. We consider this less likely given the high effect size and consistency across all 817 items.

\textbf{Token-Count Confound:} Without the token-padding control condition, we cannot empirically rule out that longer outputs (which may contain more hedging markers simply due to length) drive the correlation independently of uncertainty content.

\subsection{Honest Limitations}

\textbf{Single Model:} We validate the mechanism on GPT-3.5-turbo only. Generalization to other architectures requires explicit replication.

\textbf{Correlational Evidence:} Our findings are correlational. An intervention study manipulating hedging markers would provide stronger causal evidence.

\textbf{H-E1 Mock Mode:} The infrastructure sub-hypothesis (H-E1) verifying ECE computation ran in mock mode; actual ECE calibration improvement was not measured with real API calls. The title's reference to ``Calibration'' reflects the mechanism we verify, not end-to-end ECE improvement.

\textbf{P1 Untested:} The primary ECE comparison (CoT+confidence vs.\ single interventions) was not executed with real API calls due to resource constraints.

\textbf{Token-Padding Control Not Executed:} The planned token-padding control condition was not run. We cannot empirically rule out that calibration effects stem from token-count artifacts rather than meaningful uncertainty integration.

\textbf{Single Dataset:} Only TruthfulQA was evaluated. Cross-domain transfer is untested.

\subsection{Practical Implications}

For practitioners, our findings suggest that CoT+confidence prompting provides a zero-shot method for obtaining confidence estimates that correlate with uncertainty signals, requiring no model fine-tuning or access to internal logits. Full ECE validation remains future work.
